\chapter*{O CRP SP}

O Conselho Regional de Psicologia de São Paulo (6ª Região) é uma entidade dotada de personalidade jurídica de direito público, com autonomia administrativa e financeira, nos termos da Lei Federal n. 5.766, de 20 de dezembro de 1971. Caracteriza-se como autarquia de fiscalização do exercício e das atividades profissionais, constituindo serviço público federal, vinculado ao Conselho Federal de Psicologia (CFP) como parte do Sistema Conselhos, composto atualmente por 24 conselhos regionais e o CFP.
O CRP SP é o maior conselho regional do Sistema Conselhos, abarcando quase 1/3 da categoria profissional de psicólogas/os, com cerca de147 mil inscrições ativas de pessoas físicas e 12 mil de pessoas jurídicas. Organizado de maneira descentralizada em todo o estado de São Paulo, é dividido em uma sede administrativa, na capital, e 11 subsedes:

\begin{itemize}
    \item Alto Tietê (Mogi das Cruzes);
    \item Assis;
    \item Baixada Santista e Vale do Ribeira (Santos);
    \item Bauru;
    \item Campinas;
    \item Grande ABC (Santo André);
    \item Metropolitana (São Paulo);
    \item Ribeirão Preto;
    \item São José do Rio Preto;
    \item Sorocaba;
    \item Vale do Paraíba e Litoral Norte (Taubaté).
\end{itemize}

Além disso, também atuam no conselho 142 trabalhadoras/es efetivas/os, 11 trabalhadoras/es em cargos de livre provimento, 27 estagiárias/os, 29 prestadoras/es de serviços de limpeza e vigilância, 14 conselheiras/os efetivas/os, 15 conselheiras/os suplentes, 11 comissões gestoras (formadas por 52 gestoras/es) e mais de 380 psicólogas/os que colaboram em diversas comissões.

\section*{Missão}

Atuar com eficácia na fiscalização do exercício da profissão de psicóloga/o, com a responsabilidade de orientar, disciplinar e zelar pela fiel observância dos princípios ético-profissionais, contribuindo para o desenvolvimento da Psicologia como ciência e profissão na busca do bem comum.

\section{Visão}

Ser reconhecido pela defesa de direitos humanos e sociais , promovendo saúde e qualidade de vida para as pessoas e coletividade com responsabilidade social crítica.

\section{Valores}

Promover a liberdade, a dignidade, a igualdade e a integridade das pessoas e coletividades, contribuindo para a eliminação de quaisquer formas de negligência, crueldade e opressão, atuando para o desenvolvimento da Psicologia como campo científico de conhecimento e prática.
